\documentclass[11pt]{article}
\usepackage[margin=1in]{geometry}
\usepackage{hyperref}
\usepackage{graphicx}

% ── Paragraph spacing ────────────────────────────────────────────────────────
\usepackage{parskip}

% ── Page-break quality ───────────────────────────────────────────────────────
\widowpenalty=10000
\clubpenalty=10000
\raggedbottom

% ── Suppress automatic section numbering ─────────────────────────────────────
\setcounter{secnumdepth}{-1}

% ── Pandoc-generated table support ───────────────────────────────────────────
\usepackage{longtable}
\usepackage{booktabs}
\usepackage{array}
\usepackage{calc}
\usepackage{caption}
\newcounter{none}

% ── Math support ─────────────────────────────────────────────────────────────
\usepackage{amsmath}
\usepackage{amssymb}

% ── Make \paragraph and \subparagraph block-level ────────────────────────────
\usepackage{titlesec}
\titleformat{\paragraph}
  {\normalfont\normalsize\bfseries}{\theparagraph}{1em}{}
\titlespacing*{\paragraph}
  {0pt}{3.25ex plus 1ex minus .2ex}{1.5ex plus .2ex}
\titleformat{\subparagraph}
  {\normalfont\normalsize\bfseries}{\thesubparagraph}{1em}{}
\titlespacing*{\subparagraph}
  {0pt}{3.25ex plus 1ex minus .2ex}{1.5ex plus .2ex}

% ── Pandoc-generated tightlist support ───────────────────────────────────────
\providecommand{\tightlist}{%
  \setlength{\itemsep}{0pt}\setlength{\parskip}{0pt}}

% ── Pandoc 3.x figure sizing ────────────────────────────────────────────────
\newcommand{\pandocbounded}[1]{%
  \sbox0{#1}%
  \ifdim\wd0>\linewidth
    \resizebox{\linewidth}{!}{#1}%
  \else
    #1%
  \fi
}

% ── Unicode character fixes ──────────────────────────────────────────────────
\usepackage{newunicodechar}
\newunicodechar{✓}{$\checkmark$}
\newunicodechar{≠}{$\neq$}
\newunicodechar{≡}{$\equiv$}
\newunicodechar{≈}{$\approx$}
\newunicodechar{δ}{$\delta$}
\newunicodechar{−}{$-$}
\let\mystRightarrow\rightarrow
\renewcommand{\rightarrow}{\ensuremath{\mystRightarrow}}
% ─────────────────────────────────────────────────────────────────────────────

\hypersetup{colorlinks=true, allcolors=blue}

\title{Technical Summary: Thermodynamic Friction in Multi-Agent Resource
Contests}
\author{Keith Lostracco}
\date{}

\begin{document}
\maketitle

\subsection{Model}\label{model}

Each agent \(i\) maintains a \textbf{boundary} with integrity
\(B_i > 0\) (energy units), representing the total energetic investment
in the agent's defensive infrastructure. Maintaining boundary integrity
requires continuous work at rate
\(C_{\text{maintain},i} = \gamma_i \cdot B_i\), where \(\gamma_i > 0\)
is the entropy leakage rate. Breaching agent \(i\)'s boundary to depth
\(d \in [0,1]\) costs work
\(W_{\text{breach}}(d) = \int_0^d B_i \cdot f(s) \, ds\), where \(f\) is
a normalized resistance profile (\(\int_0^1 f(s)\,ds = 1\)), so full
breach costs \(W_{\text{breach}}(1) = B_i\).

Two agents contesting a resource \(E_j > 0\) invest energies
\(e_A, e_B \geq 0\). The probability of winning follows the Tullock
contest success function:

\[p_A(e_A, e_B) = \frac{\sigma_A \, e_A^r}{\sigma_A \, e_A^r + \sigma_B \, e_B^r}\]

where \(\sigma_i > 0\) is conflict effectiveness and \(r > 0\) is the
decisiveness parameter. The boundary integrity \(B_i\) is the
\textbf{embodied (negentropic) energy} of the boundary
(\(B_i = \Delta G_{\text{boundary}}\)), consistent with
\(W_{\text{breach}}(1) = B_i\); the damage \(\Delta B_i\) is the
embodied energy of the destroyed portion. In conflict, each agent
suffers offensive self-damage
\(\Delta B_i^{\text{off}} = \psi^{\text{off}} \cdot e_i\) and inflicts
defensive damage
\(\Delta B_j^{\text{def}} = \psi^{\text{def}} \cdot e_i\). Restoring the
boundary requires energy
\(W_{\text{repair},i} = \kappa \cdot \Delta B_i\) with repair multiplier
\(\kappa = 1 + \kappa_{\text{clear}} + \kappa_{\text{pen}} + \kappa_{\text{irr}} \geq 1\)
(\(\kappa > 1\) in any physical system, by thermodynamic
irreversibility). Because repair re-supplies the destroyed embodied
energy \(\Delta B_i\) (the leading \(1\) in \(\kappa\)), the per-agent
ledger charges contest effort plus repair only --- not a separate damage
write-off. With
\(\psi \triangleq \psi^{\text{off}} + \psi^{\text{def}}\), the
\textbf{total friction multiplier} is:

\[\Phi \triangleq 1 + \kappa(\psi^{\text{off}} + \psi^{\text{def}}) = 1 + \kappa\psi\]

The network friction coefficient \(\phi \geq 0\) parameterizes
per-transaction overhead:
\(C_{\text{overhead}}(\phi) = \phi \cdot E_{\text{transaction}}\).
Boundary violations inject friction
(\(\phi_{t+1} = \phi_t + \eta \cdot v\)), while consistent cooperation
allows friction to decay (\(\phi_{t+1} = (1 - \epsilon)\phi_t\) absent
violations).

\subsection{Results}\label{results}

\textbf{Theorem 4 (Dissipation Scaling).} In a symmetric \(N\)-agent
Tullock contest (\(r = 1\)) over resource \(E_j\):

\begin{enumerate}
\def\labelenumi{(\alph{enumi})}
\item
  Each agent's equilibrium investment:
  \(e_i^* = \frac{N-1}{N^2} \cdot E_j\).
\item
  Total contest dissipation: \(D_N = \frac{N-1}{N} \cdot E_j\).
\item
  Each agent's expected payoff: \(\Pi_i^* = \frac{E_j}{N^2}\).
\item
  Dissipation ratio: \(D_N / E_j = 1 - 1/N\).
\end{enumerate}

\textbf{Corollary 4.1 (Vanishing Individual Returns).}
\(\lim_{N \to \infty} \Pi_i^* = 0\). As the number of contestants grows,
each agent's expected return collapses to zero.

\textbf{Corollary 4.2 (Unit Dissipation Limit).}
\(\lim_{N \to \infty} D_N / E_j = 1\). In the limit of many contestants,
the entire resource value is dissipated in conflict.

\textbf{Proposition 2 (Symmetric Contest Equilibrium).} In a symmetric
Tullock contest (\(\sigma_A = \sigma_B\), \(r = 1\)) over \(E_j\):
\(e_A^* = e_B^* = E_j/4\), total dissipation \(D_2 = E_j/2\), payoffs
\(\Pi_A^* = \Pi_B^* = E_j/4\).

\textbf{Proposition 3 (Asymmetric Contest Equilibrium).} In a Tullock
contest (\(r = 1\)) with valuations \(V_A, V_B > 0\):

\[e_A^* = \frac{V_A^2 V_B}{(V_A + V_B)^2}, \qquad e_B^* = \frac{V_A V_B^2}{(V_A + V_B)^2}\]

Total dissipation: \(D_2 = V_A V_B / (V_A + V_B)\) (half the harmonic
mean).

\textbf{Theorem 5 (Net-Negative Conflict).} In a conflict contest over
resource \(E_j\) between \(N \geq 2\) symmetric agents with friction
multiplier \(\Phi\):

\begin{enumerate}
\def\labelenumi{(\alph{enumi})}
\item
  Total system loss:
  \(L_{\text{system}}^{(N)} = \Phi \cdot \frac{N-1}{N} \cdot E_j\).
\item
  Net system value:
  \(V_{\text{net}} = E_j\left(1 - \Phi \cdot \frac{N-1}{N}\right)\).
\item
  The conflict is net-negative (\(V_{\text{net}} < 0\)) if and only if
  \(\Phi > N/(N-1)\). For \(N = 2\): \(\Phi > 2\) (equivalently
  \(\kappa\psi > 1\)). For \(N = 3\): \(\Phi > 3/2\); since
  \(N/(N-1) \leq 3/2\) for all \(N \geq 3\), \(\Phi > 3/2\) suffices for
  every \(N \geq 3\). As \(N \to \infty\) the threshold falls to
  \(\Phi > 1\).
\end{enumerate}

\textbf{Corollary 5.1 (Inevitability of Net-Negative Conflict).} For any
physical system where \(\psi^{\text{off}} + \psi^{\text{def}} > 0\) and
\(\kappa > 1\), there exists a finite \(N^*\) beyond which the conflict
is net-negative.

\textbf{Corollary 5.2 (Two-Agent Threshold).} For \(N = 2\), the
conflict is net-negative when
\(\kappa\psi = \kappa(\psi^{\text{off}} + \psi^{\text{def}}) > 1\). At
\(\kappa = 2\) this requires \(\psi > 1/2\): the mild baseline
\(\psi^{\text{off}} = 0.15\), \(\psi^{\text{def}} = 0.25\)
(\(\psi = 0.4\)) gives \(\Phi = 1.8 < 2\) (net-positive), while a
destructive \(\psi^{\text{off}} = \psi^{\text{def}} = 0.3\)
(\(\psi = 0.6\)) gives \(\Phi = 2.2 > 2\) (net-negative).

\textbf{Corollary (Escalation under Physical Budgets).} For decisive
contests (\(r > N/(N-1)\)) with \(N \geq 2\) symmetric agents, the
unbounded Tullock pure NE ceases to exist
{[}@BayeKovenock-deVries1994{]}; under a finite physical energy budget
\(\bar e \leq E_j/N\), the symmetric profile \(e_i = \bar e\) is a
pure-strategy NE with \(D_N = N \bar e\). Once amplified by \(\Phi\),
the system loss exceeds the contested resource value whenever
\(\bar e > E_j/(\Phi N)\).

\textbf{Theorem 6 (Cooperation Dominance).} The cooperative regime
strictly dominates the conflict regime whenever resource collision
occurs:

\[SW_{\text{coop}} - SW_{\text{conflict}} \geq \Phi \cdot \frac{N-1}{N} \cdot E_j > 0\]

The cooperation premium is increasing in \(N\), \(\Phi\), and \(E_j\).

\textbf{Theorem 7 (Cascading Friction).} A single boundary violation of
severity \(v\) in a network with \(M\) transactions of average value
\(\bar{E}\) and friction sensitivity \(\eta\) produces a total cascading
cost of:

\[C_{\text{cascade}} = \frac{\eta \cdot v \cdot M \cdot \bar{E}}{\epsilon}\]

The ratio \(C_{\text{cascade}}/v = \eta M \bar{E} / \epsilon\) scales
linearly with network size (\(M\)) and inversely with recovery rate
(\(\epsilon\)).

\textbf{Corollary 7.1 (Friction Ratchet).} If the violation rate \(\nu\)
exceeds
\(\nu^* = \frac{\epsilon}{\eta v} \cdot \frac{E_{\text{coop net}}}{M \bar{E}}\),
the friction tax exceeds the cooperative surplus and the network
collapses.

\subsection{Notation}\label{notation}

{\def\LTcaptype{none} % do not increment counter
\begin{longtable}[]{@{}
  >{\raggedright\arraybackslash}p{(\linewidth - 2\tabcolsep) * \real{0.5000}}
  >{\raggedright\arraybackslash}p{(\linewidth - 2\tabcolsep) * \real{0.5000}}@{}}
\toprule\noalign{}
\begin{minipage}[b]{\linewidth}\raggedright
Symbol
\end{minipage} & \begin{minipage}[b]{\linewidth}\raggedright
Meaning
\end{minipage} \\
\midrule\noalign{}
\endhead
\bottomrule\noalign{}
\endlastfoot
\(B_i\) & Boundary integrity of agent \(i\) (energy units) \\
\(\gamma_i\) & Entropy leakage rate (maintenance cost coefficient) \\
\(E_j\) & Energy value of contested resource \\
\(e_i^*\) & Nash equilibrium contest investment \\
\(D_N\) & Total contest dissipation (\(N\) agents) \\
\(\psi^{\text{off}}\) & Offensive boundary-damage fraction \\
\(\psi^{\text{def}}\) & Defensive boundary-damage fraction \\
\(\psi\) & Combined boundary-damage fraction
\(\psi^{\text{off}} + \psi^{\text{def}}\) \\
\(\kappa\) & Repair multiplier (\(\kappa \geq 1\); \(\kappa > 1\) in any
physical system by the Second Law) \\
\(\Phi\) & Total friction multiplier:
\(1 + \kappa(\psi^{\text{off}} + \psi^{\text{def}})\) \\
\(\phi\) & Network friction coefficient \\
\(\eta\) & Friction sensitivity (violation-to-friction conversion) \\
\(v\) & Violation severity (energy units) \\
\(\epsilon\) & Friction recovery rate per period \\
\(M\) & Number of cooperative transactions per period \\
\(\bar{E}\) & Average cooperative transaction value \\
\end{longtable}
}

\end{document}
